RISC-V systems usually integrate general-purpose RISC-V cores with custom accelerators, memories, and communication fabrics. In such systems, performance depends not only on computation efficiency but also on how efficiently data can be transferred among these components. Typical data movement includes processor-to-accelerator configuration, memory-to-accelerator input delivery, accelerator-to-memory write-back, and accelerator-to-accelerator data exchange. In many existing systems, the RISC-V processor is responsible for orchestrating these transfers. Although this approach preserves programmability and simplifies control, it introduces extra overhead and limits scalability, especially when data transfers become frequent or when multiple accelerators must cooperate across chiplet boundaries. Therefore, data movement becomes a first-class design challenge in RISC-V systems.